% SGA 6，第 I 讲第 6—8 节及参考文献；PDF 150—169 页。
\SourcePage{150}
\SGAsection{6}{完美复形的秩}
\textbf{6.0。}设 $(X,\mathcal O_X)$ 是各环均为\Term{local-ring}{局部环}的\Term{ringed-space}{赋环空间}。众所周知，对于 $X$ 的\Term{open-subset}{开集} $U$ 上的任意\Term{finite-locally-free}{有限型局部自由层} $L$，可以赋予一个 $U$ 上的\Term{locally-constant}{局部常值}整数值函数 $\operatorname{rg}_U(L)$，称为 $L$ 的\Term{rank}{秩}。函数族 $\operatorname{rg}_U(L)$ 由下列性质刻画：
\begin{enumerate}
\item[(i)] 若 $V\subset U$ 是\Term{open-subset}{开集}的包含，则 $\operatorname{rg}_V(L|V)=\operatorname{rg}_U(L)|V$；
\item[(ii)] 对于 $\mathcal O_U$-模的每个\Term{exact-sequence}{正合列}
\[0\longrightarrow L'\longrightarrow L\longrightarrow L''\longrightarrow0,\]
其中各项均为\Term{finite-locally-free}{有限型局部自由层}，有
\[\operatorname{rg}_U(L)=\operatorname{rg}_U(L')+\operatorname{rg}_U(L'');\]
\item[(iii)] $\operatorname{rg}_X(\mathcal O_X)=1$。
\end{enumerate}
本节拟将\Term{rank}{秩}的概念推广到\Term{perfect-complex}{完美复形}。

\begin{SGAresult}{定义}{6.1}
（参见 SGA 5 VIII 2。）设 $C$ 是\Term{additive-category}{加性范畴}（相应地，\Term{triangulated-category}{三角范畴}），$F$ 是 Abel 群。若函数 $f:\operatorname{ob}(C)\to F$ 满足：对所有 $L,L',L''\in\operatorname{ob}(C)$，只要有 $L\simeq L'\oplus L''$（相应地，有\Term{distinguished-triangle}{可区别三角形} $L'\to L\to L''\to L'[1]$），就有
\[f(L)=f(L')+f(L''),\]
则称 $f$ 为从 $C$ 到 $F$ 的\Term{additive-function}{加性函数}。从 $C$ 到 $F$ 的\Term{additive-function}{加性函数}构成一个 Abel 群，记为 $\operatorname{Add}(C,F)$。
\end{SGAresult}

\SourcePage{151}
\textbf{6.2。}设 $C$ 为\Term{triangulated-category}{三角范畴}，$f$ 为从 $C$ 到 Abel 群 $F$ 的\Term{additive-function}{加性函数}。则 $f$ 在 $C$ 的底层\Term{additive-category}{加性范畴} $C^{\mathrm{ad}}$ 上诱导一个\Term{additive-function}{加性函数}，由此有包含
\[\operatorname{Add}(C,F)\subset\operatorname{Add}(C^{\mathrm{ad}},F).\]
特别地，$f(0)=0$，且若 $L\simeq M$，则 $f(L)=f(M)$。另一方面，对所有 $L\in\operatorname{ob}(C)$ 和 $n\in\mathbb Z$，有
\[f(L[n])=(-1)^nf(L).\]

\textbf{6.3。}设 $C$ 是\Term{additive-category}{加性范畴}（相应地，\Term{triangulated-category}{三角范畴}）。群 $\operatorname{Add}(C,F)$ \Term{functorial}{函子性地}依赖于 $F$。\Term{functor}{函子} $\operatorname{Add}(C,\cdot)$ 由带有\Term{additive-function}{加性映射} $\operatorname{cl}_C:\operatorname{ob}(C)\to k(C)$ 的一个 Abel 群 $k(C)$ \Term{representable-functor}{表示}（参见 SGA 5 VIII 2）。\footnote{这里不使用通常的记号 $K(C)$，因为当 $C$ 为\Term{additive-category}{加性范畴}时，$K(C)$ 已表示 $C$ 的\Term{complex}{复形}“模\Term{homotopy}{同伦}”的\Term{category}{范畴}。}（若 $C$ 为\Term{additive-category}{加性范畴}，取 $k(C)$ 为 $C$ 中\Term{object}{对象}的\Term{isomorphism}{同构}类所成\Term{monoid}{幺半群}的群\Term{completion}{完备化}；若 $C$ 为\Term{triangulated-category}{三角范畴}，取 $k(C)$ 为 $\operatorname{ob}(C)$ 生成的\Term{free-abelian-group}{自由 Abel 群}模去如下关系所得的\Term{quotient-group}{商群}：每当有\Term{distinguished-triangle}{可区别三角形} $L'\to L\to L''\to L'[1]$ 时，将 $L$ 与 $L'+L''$ 等同。）

群 $k(C)$ 相对于\Term{additive-functor}{加性}（相应地，\Term{exact}{正合}）\Term{functor}{函子}而\Term{functorial}{函子性地}依赖于 $C$。若 $u,u':C\to D$ 是\Term{isomorphism}{同构}的\Term{additive-functor}{加性}（相应地，\Term{exact}{正合}）\Term{functor}{函子}，则 $k(u)=k(u')$。

若 $C$ 是\Term{triangulated-category}{三角范畴}，则有\Term{canonical}{典范}\Term{epimorphism}{满态射} $k(C^{\mathrm{ad}})\to k(C)$（参见（6.2））。

\begin{SGAresult}{\Term{lemma}{引理}}{6.4}
设 $C$ 为\Term{additive-category}{加性范畴}。\Term{canonical}{典范}\Term{functor}{函子} $C\to K^b(C)$ 诱导一个关于 $C$ \Term{functorial}{函子性}的\Term{isomorphism}{同构}
\[k(C)\longrightarrow k(K^b(C)).\]
\end{SGAresult}
\begin{proof}
这是下列公式的直接推论（该公式由（6.2）立即得到）：
\[f(L)=\sum_i(-1)^if(L^i),\tag{6.4.1}\]
\SourcePage{152}
其中 $L\in\operatorname{ob}K^b(C)$，$f$ 是 $K^b(C)$ 上的\Term{additive-function}{加性函数}。
\end{proof}

\begin{SGAresult}{定义}{6.5}
设 $\mathscr S$ 为\Term{category}{范畴}，$\mathcal C$ 为\Term{additive-category}{加性}（相应地，\Term{triangulated-category}{三角}）$\mathscr S$-\Term{category}{范畴}（1.1.1），$F$ 为 $\mathscr S$ 上的 Abel \Term{presheaf}{预层}。若以 $X\in\operatorname{ob}\mathscr S$ 为指标的函数族 $f_X:\operatorname{ob}(\mathcal C_X)\to F(X)$ 满足下列条件，则称其为从 $\mathcal C$ 到 $F$ 的\Term{additive-s-function}{$\mathscr S$-加性函数}：
\begin{enumerate}
\item[(i)] 对每个 $X\in\operatorname{ob}\mathcal C_X$，$f_X$ 是\Term{additive-function}{加性函数}（6.1）；
\item[(ii)] 对 $\mathscr S$ 的每个\Term{arrow}{箭头} $u:X\to Y$，下图\Term{commutative-diagram}{交换}：
\[\begin{tikzcd}
\operatorname{ob}\mathcal C_Y\arrow[r,"f_Y"]\arrow[d,"u^*"']&F(Y)\arrow[d,"F(u)"]\\
\operatorname{ob}\mathcal C_X\arrow[r,"f_X"']&F(X).
\end{tikzcd}\]
\end{enumerate}
从 $\mathcal C$ 到 $F$ 的\Term{additive-s-function}{$\mathscr S$-加性函数}构成一个 Abel 群，记作 $\operatorname{Add}_{\mathscr S}(\mathcal C,F)$。
\end{SGAresult}

\textbf{注记 6.6。}a）以 $\mathcal C^{\mathrm{is}}$ 表示如下 $\mathscr S$-\Term{fibered-category}{纤维范畴}：其在 $X\in\operatorname{ob}\mathscr S$ 处的\Term{fiber}{纤维}，是与 $\mathcal C_X$ 有相同\Term{object}{对象}，而\Term{arrow}{箭头}仅取 $\mathcal C_X$ 的\Term{isomorphism}{同构}的\Term{category}{范畴}。另一方面，以 $\mathcal F$ 表示如下 $\mathscr S$-\Term{fibered-category}{纤维范畴}：其在 $X\in\operatorname{ob}\mathscr S$ 处的\Term{fiber}{纤维}为\Term{discrete-category}{离散范畴}（即仅有恒等\Term{arrow}{箭头}）$F(X)$，而沿 $u\in\operatorname{Fl}(\mathscr S)$ 的\Term{inverse-image}{逆像}\Term{functor}{函子}为 $F(u)$。于是，从 $\mathcal C$ 到 $F$ 的一个\Term{additive-s-function}{$\mathscr S$-加性函数}，正是从 $\mathcal C^{\mathrm{is}}$ 到 $\mathcal F$ 的一个\Term{cartesian}{笛卡尔} $\mathscr S$-\Term{functor}{函子}，它在每个\Term{fiber}{纤维}上诱导（6.1）意义的\Term{additive-function}{加性函数}。以下解释\Term{additive-function}{加性}的方式将很方便。为确定起见，设 $\mathcal C$ 为\Term{triangulated-category}{三角}的。由从 $\mathcal C^{\mathrm{is}}$（相应地，$\operatorname{Tr}(\mathcal C)^{\mathrm{is}}$）到 $\mathcal F$ 的\Term{cartesian}{笛卡尔} $\mathscr S$-\Term{functor}{函子}组成的\Term{category}{范畴}
$\operatorname{Cart}_{\mathscr S}(\mathcal C^{\mathrm{is}},\mathcal F)$（相应地，$\operatorname{Cart}_{\mathscr S}(\operatorname{Tr}(\mathcal C)^{\mathrm{is}},\mathcal F)$）（1.1.2），是一个 Abel 群所对应的\Term{discrete-category}{离散范畴}，
\SourcePage{153}
且有\Term{homomorphism}{同态}
\[\rho:\operatorname{Cart}_{\mathscr S}(\mathcal C^{\mathrm{is}},\mathcal F)\longrightarrow\operatorname{Cart}_{\mathscr S}(\operatorname{Tr}(\mathcal C)^{\mathrm{is}},\mathcal F),\]
其定义为：对于 $f\in\operatorname{Cart}_{\mathscr S}(\mathcal C^{\mathrm{is}},\mathcal F)$ 和
$T=(L'\to L\to L''\to L'[1])\in\operatorname{ob}\operatorname{Tr}(\mathcal C)$，令
\[\rho(f)(T)=f(L)-f(L')-f(L'').\]
称 $f\in\operatorname{Cart}_{\mathscr S}(\mathcal C^{\mathrm{is}},\mathcal F)$ 为\Term{additive-function}{加性}的，意即 $\rho(f)=0$；换言之，有\Term{canonical}{典范}且\Term{functorial}{函子性}的\Term{exact-sequence}{正合列}
\[0\longrightarrow\operatorname{Add}_{\mathscr S}(\mathcal C,F)\longrightarrow\operatorname{Cart}_{\mathscr S}(\mathcal C^{\mathrm{is}},\mathcal F)\xrightarrow{\rho}\operatorname{Cart}_{\mathscr S}(\operatorname{Tr}(\mathcal C)^{\mathrm{is}},\mathcal F).\tag{6.6.1}\]

b）设 $u:X\to Y$ 为 $\mathscr S$ 的\Term{arrow}{箭头}。所有“沿 $u$ 的\Term{inverse-image}{逆像}”\Term{functor}{函子}彼此\Term{isomorphism}{同构}，故由（6.3）定义一个唯一的\Term{homomorphism}{同态} $k(Y)\to k(X)$。于是 $X\mapsto k(X)$ 定义 $\mathscr S$ 上的一个 Abel \Term{presheaf}{预层}，记为 $k(\mathcal C)$。从 $\mathcal C$ 到 $F$ 的一个\Term{additive-s-function}{$\mathscr S$-加性函数}，就简单地解释为\Term{presheaf}{预层}\Term{homomorphism}{同态} $k(\mathcal C)\to F$；换言之，有
\[\operatorname{Add}_{\mathscr S}(\mathcal C,F)=\operatorname{Hom}(k(\mathcal C),F).\]
当然，\Term{presheaf}{预层} $k(\mathcal C)$ 相对于\Term{additive-functor}{加性}（相应地，\Term{exact}{正合}）$\mathscr S$-\Term{functor}{函子}而\Term{functorial}{函子性地}依赖于 $\mathcal C$。

\begin{SGAresult}{命题}{6.7}
设 $\mathscr S$ 为\Term{site}{景}，$\mathcal C$ 为\Term{flat}{平坦} Abel $\mathscr S$-\Term{category}{范畴}，$\mathcal C_0$ 为 $\mathcal C$ 的\Term{additive-category}{加性}严格\Term{full-subcategory}{满子范畴}，满足假设（4.0）。则由 $\mathcal C_0$ 到 $\operatorname{Parf}(\mathcal C,\mathcal C_0)$ 的包含所诱导的\Term{presheaf}{预层}\Term{morphism}{态射}
\[k(\mathcal C_0)\longrightarrow k(\operatorname{Parf}(\mathcal C,\mathcal C_0))\qquad\text{（参见（4.9））}\]
在\Term{associated-sheaf}{伴随层}上定义一个\Term{isomorphism}{同构}。换言之，对于 $\mathscr S$ 上任意 Abel \Term{sheaf}{层} $F$，
\SourcePage{154}
\Term{restriction}{限制}\Term{homomorphism}{同态}
\[\operatorname{Add}_{\mathscr S}(\operatorname{Parf}(\mathcal C,\mathcal C_0),F)\longrightarrow\operatorname{Add}_{\mathscr S}(\mathcal C_0,F)\]
是\Term{isomorphism}{同构}。
\end{SGAresult}
\begin{proof}
设 $F$ 为 $\mathscr S$ 上的 Abel \Term{sheaf}{层}，证明\Term{arrow}{箭头}
\[\operatorname{Add}_{\mathscr S}(\operatorname{Parf}(\mathcal C,\mathcal C_0),F)\longrightarrow\operatorname{Add}_{\mathscr S}(\mathcal C_0,F)\]
为\Term{isomorphism}{同构}。它\Term{resolution}{分解}为
\[\operatorname{Add}_{\mathscr S}(\operatorname{Parf}(\mathcal C,\mathcal C_0),F)\xrightarrow{a}\operatorname{Add}_{\mathscr S}(K^b(\mathcal C_0),F)\xrightarrow{b}\operatorname{Add}_{\mathscr S}(\mathcal C_0,F).\]
由（6.4），\Term{arrow}{箭头} $b$ 为\Term{isomorphism}{同构}。故归结为证明 $a$ 是\Term{isomorphism}{同构}。

考虑下列\Term{commutative-diagram}{交换图}，其中各行\Term{exact}{正合}（6.6.1），竖\Term{arrow}{箭头}为“\Term{restriction}{限制}”\Term{arrow}{箭头}：
\[\resizebox{\textwidth}{!}{\begin{tikzcd}[ampersand replacement=\&,column sep=small]
0\arrow[r]\&\operatorname{Add}_{\mathscr S}(\operatorname{Parf}(\mathcal C),F)\arrow[r]\arrow[d,"a"]\&\operatorname{Cart}_{\mathscr S}(\operatorname{Parf}(\mathcal C)^{\mathrm{is}},\mathcal F)\arrow[r]\arrow[d]\&\operatorname{Cart}_{\mathscr S}(\operatorname{Tr}(\operatorname{Parf}(\mathcal C))^{\mathrm{is}},\mathcal F)\arrow[d]\\
0\arrow[r]\&\operatorname{Add}_{\mathscr S}(K^b(\mathcal C_0),F)\arrow[r]\&\operatorname{Cart}_{\mathscr S}(K^b(\mathcal C_0)^{\mathrm{is}},\mathcal F)\arrow[r]\&\operatorname{Cart}_{\mathscr S}(\operatorname{Tr}(K^b(\mathcal C_0))^{\mathrm{is}},\mathcal F).
\end{tikzcd}}\]
因为 $F$ 是\Term{sheaf}{层}，$\mathcal F$ 是\Term{stack}{叠}（G II 2.2.4），故由（4.10），右侧两个竖\Term{arrow}{箭头}均为\Term{isomorphism}{同构}。所以 $a$ 亦如此，证毕。
\end{proof}

\textbf{注记 6.7.1。}（参见（4.10.1）。）在（6.7）中，可以将 $\operatorname{Parf}(\mathcal C,\mathcal C_0)$ 替换为 $D^*(\mathcal C)_{\mathcal C_0\text{-}\mathrm{parf}}$，其中 $*= -$ 或 $b$。

\SourcePage{155}
\textbf{6.8。}设 $(\mathscr S,A)$ 为\Term{ringed-topos}{赋环拓扑斯}，$\mathcal C$ 和 $\mathcal C_0$ 为（4.8）中所考虑的 $\mathscr S$-\Term{fibered-category}{纤维范畴}。有 Abel \Term{presheaf}{预层}\Term{homomorphism}{同态}
\[k(\mathcal C_0)\xrightarrow{\alpha}k(D^-({}_A\mathscr S)_{\mathrm{parf}})\xrightarrow{\beta}k(\operatorname{Parf}({}_A\mathscr S)),\]
由（6.7），它们在\Term{associated-sheaf}{伴随层}上诱导\Term{isomorphism}{同构}。\Term{homomorphism}{同态} $\alpha$ 与\Term{inverse-image}{逆像}相容，即若
$u:(\mathscr S',A')\to(\mathscr S,A)$ 是\Term{ringed-topos}{赋环拓扑斯}的\Term{morphism}{态射}，则有下列\Term{commutative-diagram}{交换图}，其中竖\Term{arrow}{箭头}由\Term{functor}{函子} $\mathrm Lu^*$ 诱导（参见（4.19.1））：
\[\begin{tikzcd}
k(\mathcal C_0)\arrow[r]\arrow[d]&k(D^-({}_A\mathscr S)_{\mathrm{parf}})\arrow[d]\\
k(\mathcal C'_0)\arrow[r]&k(D^-({}_{A'}\mathscr S')_{\mathrm{parf}}).
\end{tikzcd}\]
（$\mathcal C'_0$ 的定义与 $\mathcal C_0$ 类似。）若 $A$ 交换，则\Term{tensor-product}{张量积} $\mathbin{\otimes^{\mathrm L}}$ 在 $k(\mathcal C_0)$ 和 $k(D^-({}_A\mathscr S)_{\mathrm{parf}})$ 上赋予环\Term{presheaf}{预层}结构（参见（4.19.1）），且 $\alpha$ 是环\Term{presheaf}{预层}\Term{homomorphism}{同态}。

\textbf{6.9。}在（6.8）的情形，假定 $(\mathscr S,A)$ 是\Term{locally-ringed-topos}{局部赋环拓扑斯}（参见（2.15.1））。则由（2.15.1），对于每个 $X\in\operatorname{ob}\mathscr S$，$\mathcal C_{0X}$ 是\Term{finite-locally-free}{有限型局部自由} $A_X$-模的\Term{category}{范畴}。以 $\mathcal C_1$ 表示由 $X\mapsto$（\Term{finite-free}{有限型自由} $A_X$-模的\Term{category}{范畴}）定义的 $\mathscr S$-\Term{fibered-category}{纤维范畴}。显然 $\mathcal C_1$ 满足假设（4.0）。$D(\mathscr S)$ 的\Term{object}{对象}为\Term{perfect-complex}{完美的}，当且仅当它相对于 $\mathcal C_1$ 为\Term{perfect-complex}{完美的}。有 Abel \Term{presheaf}{预层}\Term{homomorphism}{同态}
\[k(\mathcal C_1)\longrightarrow k(\mathcal C_0)\longrightarrow k(D^-(\mathscr S)_{\mathrm{parf}})\longrightarrow k(\operatorname{Parf}(\mathscr S))\]
\SourcePage{156}
（左侧两个是环\Term{presheaf}{预层}\Term{homomorphism}{同态}），由（6.7）它们在\Term{associated-sheaf}{伴随层}上诱导\Term{isomorphism}{同构}。下面证明这些\Term{associated-sheaf}{伴随层}实际上\Term{canonical}{典范}\Term{isomorphism}{同构}于\Term{constant-sheaf}{常值层} $\mathbb Z_{\mathscr S}$。

\begin{SGAresult}{\Term{lemma}{引理}}{6.10}
设 $(\mathscr S,A)$ 为\Term{locally-ringed-topos}{局部赋环拓扑斯}，$U$ 为 $\mathscr S$ 中异于空\Term{object}{对象}的一个\Term{object}{对象}，$p,q\in\mathbb N$。若 $A_U^p\simeq A_U^q$，则 $p=q$。
\end{SGAresult}
\begin{proof}
用与（2.15.1）证明类似的论证，归结到 $(\mathscr S,A)$ 为由一个\Term{prescheme}{预概形}定义的\Term{locally-ringed-topos}{局部赋环拓扑斯}的情形；此时断言显然（在 $U$ 中取一点）。
\end{proof}

\begin{SGAresult}{定义}{6.11}
由下式定义的从 $\mathcal C_1$ 到 $\mathbb Z_{\mathscr S}$ 的\Term{additive-s-function}{$\mathscr S$-加性函数}称为\Term{rank}{秩}，记作 $\operatorname{rg}$：
\[\begin{aligned}
\operatorname{rg}(L)&=p&&\text{若 }L\simeq A_X^p,\quad X\ne\varnothing,\\
\operatorname{rg}(L)&=0&&\text{若 }L\in\operatorname{ob}\mathcal C_{1X},\quad X=\varnothing.
\end{aligned}\]
由（6.7），\Term{rank}{秩}函数唯一延拓为从 $\operatorname{Parf}(\mathscr S)$ 到 $\mathbb Z_{\mathscr S}$ 的\Term{additive-s-function}{$\mathscr S$-加性函数}（我们仍称其为\Term{rank}{秩}，记作 $\operatorname{rg}$）。
\end{SGAresult}

\begin{SGAresult}{命题}{6.12}
a）对于 $L,M\in\operatorname{ob}D^-(X)_{\mathrm{parf}}$（$X\in\operatorname{ob}\mathscr S$），有
\[\operatorname{rg}(L\otimes^{\mathrm L}M)=\operatorname{rg}(L)\operatorname{rg}(M).\]
b）设 $u:(\mathscr S',A')\to(\mathscr S,A)$ 是\Term{locally-ringed-topos}{局部赋环拓扑斯}的\Term{morphism}{态射}。对于 $M\in\operatorname{ob}D^-(\mathscr S)_{\mathrm{parf}}$，有
\[\operatorname{rg}(\mathrm Lu^*(M))=u^*(\operatorname{rg}(M)).\]
换言之，\Term{rank}{秩}\Term{homomorphism}{同态}
\[\operatorname{rg}:k(D^-(\mathscr S)_{\mathrm{parf}})\longrightarrow\mathbb Z_{\mathscr S}\]
是与\Term{inverse-image}{逆像}相容的环\Term{presheaf}{预层}\Term{homomorphism}{同态}。
\end{SGAresult}
\SourcePage{157}
\begin{proof}
只须证明\Term{rank}{秩}在 $k(D^-(\mathscr S)_{\mathrm{parf}})$ 的\Term{associated-sheaf}{伴随层}上诱导的\Term{homomorphism}{同态}是与\Term{inverse-image}{逆像}相容的环\Term{sheaf}{层}\Term{homomorphism}{同态}。由（6.9），可以将 $k(D^-(\mathscr S)_{\mathrm{parf}})$ 替换为 $k(\mathcal C_1)$。于是归结为对\Term{finite-free}{有限型自由}的 $L$ 和 $M$ 验证 a）与 b），这是显然的。
\end{proof}

\textbf{6.13。}$(\mathscr S,A)$ 仍表示\Term{locally-ringed-topos}{局部赋环拓扑斯}，设 $X\in\operatorname{ob}\mathscr S$，$\underline n\in H^0(X,\mathbb Z_{\mathscr S})$。如下定义 $\operatorname{Parf}(X)$ 的一个\Term{object}{对象}，记作 $A_X^{\underline n}$。若 $\underline n$ 为\Term{constant}{常值}，值为整数 $n\geq0$（相应地，$n<0$），则置
\[A_X^{\underline n}=A_X^n\qquad\text{（相应地，}A_X^{\underline n}=A_X^{-n}[1]\text{）；}\]
\Term{generic}{一般}情形，将 $X$ \Term{resolution}{分解}为不交的部分 $X_i$，使 $\underline n|X_i$ 为\Term{constant}{常值}，并将 $A_X^{\underline n}$ 定义为 $\operatorname{Parf}(X)$ 中在每个 $X_i$ 上诱导 $A_{X_i}^{\underline n|X_i}$ 的唯一\Term{object}{对象}。（更\Term{generic}{一般}地，对于 $L\in\operatorname{ob}\operatorname{Parf}(X)$，可令 $L^{\underline n}=A_X^{\underline n}\otimes^{\mathrm L}L$，从而定义 $L^{\underline n}$。）由 $i_X(\underline n)=\operatorname{cl}(A_X^{\underline n})$ 定义的映射
\[i_X:H^0(X,\mathbb Z_{\mathscr S})\longrightarrow k(\operatorname{Parf}(X))\]
是\Term{ring-homomorphism}{环同态}（且 $i_X(1)=1$），并在 $X$ 变动时定义环\Term{presheaf}{预层}\Term{homomorphism}{同态}
\[i:\mathbb Z_{\mathscr S}\longrightarrow k(\operatorname{Parf}(\mathscr S)).\]

\begin{SGAresult}{命题}{6.14}
在（6.13）的假设下，有
\[\operatorname{rg}\,i=\operatorname{Id}(\mathbb Z_{\mathscr S})\]
（换言之，借助 $i$ 视为 $\mathbb Z_{\mathscr S}$-\Term{algebra}{代数}\Term{presheaf}{预层}的 $k(\operatorname{Parf}(\mathscr S))$，通过\Term{rank}{秩}\Term{augmentation}{增广}到 $\mathbb Z_{\mathscr S}$）。此外，\Term{rank}{秩}
\[\operatorname{rg}:k(\operatorname{Parf}(\mathscr S))\longrightarrow\mathbb Z_{\mathscr S}\]
在\Term{associated-sheaf}{伴随层}上定义一个\Term{isomorphism}{同构}。
\end{SGAresult}
\SourcePage{158}
\begin{proof}
第一断言由下列关系立即得到：对于 $L\in\operatorname{ob}\operatorname{Parf}(\mathscr S)$，
\[\operatorname{cl}(L[1])=-\operatorname{cl}(L),\qquad\operatorname{cl}(L^n)=n\operatorname{cl}(L).\]
证明第二断言。首先注意到，$i$ 显然经过 $k(\mathcal C_0)$ \Term{resolution}{分解}。因为 $\operatorname{rg}\,i$ 为恒等，由（6.9）归结为证明\Term{homomorphism}{同态}
\[i\operatorname{rg}:k(\mathcal C_0)\longrightarrow k(\mathcal C_0)\]
在\Term{associated-sheaf}{伴随层}上诱导\Term{isomorphism}{同构}；即对 $L\in\operatorname{ob}\mathcal C_0$，$i\operatorname{rg}(L)$ \Term{local-scheme}{局部}\Term{isomorphism}{同构}于 $L$，这是立即可见的。
\end{proof}
由（6.6 b））得到：
\begin{SGAresult}{推论}{6.15}
在（6.13）的假设下，对于 $\mathscr S$ 上任意 Abel \Term{sheaf}{层} $F$，从 $\operatorname{Parf}(\mathscr S)$ 到 $F$ 的\Term{additive-s-function}{$\mathscr S$-加性函数}群\Term{canonical}{典范}地等同于 $H^0(S,F)$（$S$ 是 $\mathscr S$ 的\Term{final-object}{终对象}）；\Term{rank}{秩}函数对应于 $F=\mathbb Z_{\mathscr S}$ 及 $\mathbb Z_{\mathscr S}$ 的\Term{unit}{单位}\Term{section}{截面}。
\end{SGAresult}

\SGAsection{7}{完美复形的对偶性}
\textbf{7.0。}本节中，$(\mathscr S,A)$ 表示一个环均为\Term{commutative-ring}{交换环}的\Term{ringed-topos}{赋环拓扑斯}，以 $S$ 表示 $\mathscr S$ 的\Term{final-object}{终对象}。我们拟将\Term{finite-locally-free}{有限型局部自由层}的熟知\Term{duality}{对偶性}概念推广到\Term{perfect-complex}{完美复形}。

\begin{SGAresult}{命题}{7.1}
设 $E\in\operatorname{ob}D(S)_{\mathrm{parf}}$，$F\in\operatorname{ob}D^+(S)$。

a）若 $F$ 的\Term{cohomology}{上同调}\Term{local-scheme}{局部}\Term{bounded}{有界}，则 $\operatorname{R\mathcal Hom}(E,F)$ 亦如此。若 $F\in\operatorname{ob}D(S)$ 是\Term{pseudo-coherent}{伪凝聚}的（相应地，\Term{local-scheme}{局部}具有\Term{finite}{有限}\Term{tor-dimension}{Tor 维数}；相应地，\Term{perfect-complex}{完美}），则 $\operatorname{R\mathcal Hom}(E,F)$ 亦如此。

\SourcePage{159}
b）再假定 $E$ 具有\Term{finite}{有限}\Term{tor-dimension}{Tor 维数}（若 $\mathscr S$ 为\Term{finite-type}{有限型}，此条件即满足）。则若 $F$ 的\Term{cohomology}{上同调}\Term{bounded}{有界}（相应地，具有\Term{finite}{有限}\Term{tor-dimension}{Tor 维数}），$\operatorname{R\mathcal Hom}(E,F)$ 亦如此。更精确地，设整数 $a,b,m,n$ 满足
\[\operatorname{tor.amp}(E)\subset[a,b]\]
（或者，\Term{equivalence}{等价}地，由（5.8），$\operatorname{parf.amp}(E)\subset[a,b]$），且 $i\notin[m,n]$ 时 $H^i(F)=0$（相应地，$\operatorname{tor.amp}(F)\subset[m,n]$）。则当 $i\notin[m-b,n-a]$ 时，有
\[H^i(\operatorname{R\mathcal Hom}(E,F))=0\]
（相应地，$\operatorname{tor.amp}(\operatorname{R\mathcal Hom}(E,F))\subset[m-b,n-a]$）。
\end{SGAresult}
\begin{proof}
断言 a）是\Term{local-scheme}{局部}的，故可设 $E$ 为\Term{strictly-perfect}{严格完美}的；对 $E$ 作 dévissage，归结为 $E=A$，即可得证。
对于断言 b），可设 $E$ 为\Term{strictly-perfect}{严格完美}的，且 $i\notin[a,b]$ 时 $E^i=0$；并且 $i\notin[m,n]$ 时 $F^i=0$（相应地，每个 $F^i$ 均\Term{flat}{平坦}，且 $i\notin[m,n]$ 时 $F^i=0$）。此时只须应用下列\Term{lemma}{引理}。
\end{proof}

\begin{SGAresult}{\Term{lemma}{引理}}{7.1.1}
设 $E\in\operatorname{ob}D^-(S)$，$F\in\operatorname{ob}D^+(S)$。若 $E$ 为\Term{strictly-pseudo-coherent}{严格伪凝聚}的（即（2.1），各\Term{set-component}{分量}均为\Term{finite-free}{有限型自由模}的\Term{direct-summand}{直和因子}），则 $D^+(S)$ 的\Term{canonical}{典范}\Term{arrow}{箭头}
\[\mathcal Hom^\bullet(E,F)\longrightarrow\operatorname{R\mathcal Hom}(E,F)\]
是\Term{isomorphism}{同构}。
\end{SGAresult}
\begin{proof}
\Term{spectral-sequence}{谱序列}
\[E_2^{pq}=H^p(\mathcal Ext^q(E,F)^\bullet)\Longrightarrow\mathcal Ext^*(E,F),\]
其中
\[\mathcal Ext^q(E,F)^i=\prod_{i_2-i_1=i}\mathcal Ext^q(E^{i_1},F^{i_2}),\]
\Term{degeneration}{退化}，因为对所有 $q>0$ 有 $\mathcal Ext^q(E,F)^\bullet=0$（1.3.1）。由此得到\Term{isomorphism}{同构}
\SourcePage{160}
\[E_2^{p0}=H^p(\mathcal Hom^\bullet(E,F))\xrightarrow{\sim}\mathcal Ext^p(E,F)=H^p(\operatorname{R\mathcal Hom}(E,F)),\]
从而证明了断言。
\end{proof}

\begin{SGAresult}{命题}{7.1.2}
设 $f:(\mathscr S',A')\to(\mathscr S,A)$ 是\Term{ringed-topos}{赋环拓扑斯}的\Term{morphism}{态射}。设 $E\in\operatorname{ob}D^-(S)$ 和 $F\in\operatorname{ob}D^b(S)$ 满足
\[\operatorname{R\mathcal Hom}(E,F)\in\operatorname{ob}D^b(S),\qquad\mathrm Lf^*(F)\in\operatorname{ob}D^b(S').\]
则存在\Term{canonical}{典范}且\Term{functorial}{函子性}的\Term{arrow}{箭头}
\[\mathrm Lf^*\operatorname{R\mathcal Hom}(E,F)\longrightarrow\operatorname{R\mathcal Hom}(\mathrm Lf^*E,\mathrm Lf^*F).\]
若 $E$ 为\Term{perfect-complex}{完美的}，则此\Term{arrow}{箭头}为\Term{isomorphism}{同构}。
\end{SGAresult}
\begin{proof}
由\Term{functor}{函子} $\mathrm Lf^*$ 与 $\mathrm Rf_*$ 的\Term{adjunction}{伴随性}，只须定义\Term{arrow}{箭头}
\[\operatorname{R\mathcal Hom}(E,F)\longrightarrow\mathrm Rf_*\operatorname{R\mathcal Hom}(\mathrm Lf^*E,\mathrm Lf^*F).\]
为此，使用\Term{adjunction}{伴随}\Term{arrow}{箭头} $F\to\mathrm Rf_*\mathrm Lf^*F$ 及\Term{duality}{对偶性}\Term{isomorphism}{同构}
\[\mathrm Rf_*\operatorname{R\mathcal Hom}(\mathrm Lf^*E,\mathrm Lf^*F)\xrightarrow{\sim}\operatorname{R\mathcal Hom}(E,\mathrm Rf_*\mathrm Lf^*F).\]
为说明当 $E$ 为\Term{perfect-complex}{完美}时（7.1.2）的\Term{arrow}{箭头}是\Term{isomorphism}{同构}，经\Term{localization}{局部化}和 dévissage，立即归结到 $E=A$；此时结论是恒真的。
\end{proof}

\begin{SGAresult}{命题}{7.2}
对于 $E\in\operatorname{ob}D(S)$ 和 $F\in\operatorname{ob}D^+(S)$，存在\Term{canonical}{典范}且\Term{functorial}{函子性}的\Term{arrow}{箭头}（参见 SGA 5 I 1.6）
\[E\longrightarrow\operatorname{R\mathcal Hom}(\operatorname{R\mathcal Hom}(E,F),F).\]
若 $E\in\operatorname{ob}D(S)_{\mathrm{parf}}$ 且 $F=A$，则此\Term{arrow}{箭头}为\Term{isomorphism}{同构}。
\end{SGAresult}
\SourcePage{161}
\begin{proof}
对 $F$ 作\Term{injective}{内射}\Term{resolution}{消解}，并使用 $K(S)$ 中的\Term{canonical}{典范}\Term{arrow}{箭头}
\[E\longrightarrow\mathcal Hom^\bullet(\mathcal Hom^\bullet(E,F),F),\]
即定义所需\Term{arrow}{箭头}。为说明当 $E$ 为\Term{perfect-complex}{完美}且 $F=A$ 时该\Term{arrow}{箭头}是\Term{isomorphism}{同构}，经\Term{localization}{局部化}与 dévissage，归结到 $E=A$，此时结论显然。
\end{proof}

\begin{SGAresult}{定义}{7.3}
设 $E\in\operatorname{ob}D(S)$。置
\[\operatorname{R\mathcal Hom}(E,A)=E^\vee,\]
称 $E^\vee$ 为 $E$ 的\Term{dual}{对偶}。
\end{SGAresult}
由（7.1）与（7.2），若 $E\in\operatorname{ob}D(S)$ 是\Term{perfect-complex}{完美的}（相应地，\Term{perfect-complex}{完美}且具有\Term{finite}{有限}\Term{tor-dimension}{Tor 维数}），则 $E^\vee$ 亦如此，且\Term{canonical}{典范}\Term{arrow}{箭头} $E\to E^{\vee\vee}$ 是\Term{isomorphism}{同构}。换言之，\Term{functor}{函子} $\operatorname{R\mathcal Hom}(\,\cdot\,,A)$ 定义\Term{category}{范畴} $D(S)_{\mathrm{parf}}$（相应地，$D(S)_{\underline{\mathrm{parf}}}$（5.8.4 b）））到自身的一个\Term{anti-equivalence}{反自等价}。

\begin{SGAresult}{\Term{lemma}{引理}}{7.4}
（Cartan \Term{isomorphism}{同构}；参见 [H] II 5.15。）对于 $E\in\operatorname{ob}D(S)$、$F\in\operatorname{ob}D^-(S)$、$G\in\operatorname{ob}D^+(S)$，存在\Term{canonical}{典范}且\Term{functorial}{函子性}的\Term{isomorphism}{同构}：
\[\begin{aligned}
\operatorname{R\mathcal Hom}(E\otimes^{\mathrm L}F,G)&\xrightarrow{\sim}\operatorname{R\mathcal Hom}(E,\operatorname{R\mathcal Hom}(F,G)),\\
\operatorname{RHom}(E\otimes^{\mathrm L}F,G)&\xrightarrow{\sim}\operatorname{RHom}(E,\operatorname{R\mathcal Hom}(F,G)),\\
\operatorname{Hom}(E\otimes^{\mathrm L}F,G)&\xrightarrow{\sim}\operatorname{Hom}(E,\operatorname{R\mathcal Hom}(F,G)).
\end{aligned}\]
\end{SGAresult}
\SourcePage{162}
\begin{proof}
这是标准事实：对 $G$ 作\Term{injective}{内射}\Term{resolution}{消解}，对 $F$ 作\Term{flat}{平坦}\Term{resolution}{消解}。
\end{proof}

\begin{SGAresult}{\Term{lemma}{引理}}{7.5}
对于 $E\in\operatorname{ob}D^-(S)$、$F\in\operatorname{ob}D^+(S)$，存在\Term{canonical}{典范}且\Term{functorial}{函子性}的\Term{homomorphism}{同态}
\[E\otimes^{\mathrm L}\operatorname{R\mathcal Hom}(E,F)\longrightarrow F.\]
\end{SGAresult}

\begin{SGAresult}{命题}{7.6}
（参见 [H] II 5.13。）设 $E\in\operatorname{ob}D(S)$、$F\in\operatorname{ob}D^+(S)$、$G\in\operatorname{ob}D^+(S)$。在下列任一假设下：
\begin{enumerate}
\item[(i)] $G$ 具有\Term{finite}{有限}\Term{tor-dimension}{Tor 维数}；
\item[(ii)] $E\in\operatorname{ob}D^-(S)$，$F$ 具有\Term{finite}{有限}\Term{tor-dimension}{Tor 维数}，且 $G\in\operatorname{ob}D^b(S)$，
\end{enumerate}
存在\Term{canonical}{典范}且\Term{functorial}{函子性}的\Term{homomorphism}{同态}
\[\operatorname{R\mathcal Hom}(E,F)\otimes^{\mathrm L}G\longrightarrow\operatorname{R\mathcal Hom}(E,F\otimes^{\mathrm L}G).\]
当 $E$ 为\Term{perfect-complex}{完美的}时，此\Term{homomorphism}{同态}为\Term{isomorphism}{同构}。
\end{SGAresult}
\begin{proof}
先定义\Term{arrow}{箭头}。在假设（i）下，只须对 $F$ 作\Term{injective}{内射}\Term{resolution}{消解}、对 $G$ 作\Term{flat}{平坦}\Term{resolution}{消解}。在假设（ii）下，由（7.4），\Term{equivalence}{等价}于定义\Term{arrow}{箭头}
\[E\otimes^{\mathrm L}\operatorname{R\mathcal Hom}(E,F)\otimes^{\mathrm L}G\longrightarrow F\otimes^{\mathrm L}G.\]
对（7.5）的\Term{arrow}{箭头}应用\Term{functor}{函子} $\mathbin{\otimes^{\mathrm L}}G$ 即可。

为说明当 $E$ 为\Term{perfect-complex}{完美}时此\Term{arrow}{箭头}为\Term{isomorphism}{同构}，使用标准方法（参见（7.2）的证明）：经\Term{localization}{局部化}与 dévissage，归结为 $E=A$。
\end{proof}

\SourcePage{163}
\begin{SGAresult}{推论}{7.7}
对于 $E\in\operatorname{ob}D(S)_{\underline{\mathrm{parf}}}$ 与 $F\in\operatorname{ob}D^+(S)$，存在\Term{functorial}{函子性}\Term{isomorphism}{同构}
\[E^\vee\otimes^{\mathrm L}F\xrightarrow{\sim}\operatorname{R\mathcal Hom}(E,F),\]
其中使用（7.3）的记号。
\end{SGAresult}
\begin{proof}
应用（7.6），其中 $F$ 和 $G$ 分别替换为 $A$ 和 $F$；并注意到由（7.1），$E^\vee\in\operatorname{ob}D(S)_{\underline{\mathrm{parf}}}$。
\end{proof}

\begin{SGAresult}{推论}{7.8}
设 $E,F\in\operatorname{ob}D(S)_{\underline{\mathrm{parf}}}$。存在自然\Term{arrow}{箭头}
\[\operatorname{R\mathcal Hom}(E,F)\otimes^{\mathrm L}\operatorname{R\mathcal Hom}(F,E)\longrightarrow A,\]
它是一个“\Term{perfect-duality}{完美对偶}”，即借助（7.4），将两个\Term{complex}{复形} $\operatorname{R\mathcal Hom}(E,F)$、$\operatorname{R\mathcal Hom}(F,E)$ 中的每一个等同于另一个的\Term{dual}{对偶}（按（7.3）的意义）。
\end{SGAresult}
\begin{proof}
由（7.5），有自然\Term{pairing}{配对} $E^\vee\otimes^{\mathrm L}E\to A$、$F^\vee\otimes^{\mathrm L}F\to A$，从而通过\Term{tensor-product}{张量积}得到自然\Term{pairing}{配对}
\[E^\vee\otimes^{\mathrm L}F\otimes^{\mathrm L}F^\vee\otimes^{\mathrm L}E\longrightarrow A.\]
考虑（7.7），这就给出所宣称的 $\operatorname{R\mathcal Hom}$ 上的\Term{pairing}{配对}。
证明它是\Term{perfect-duality}{完美对偶}。有\Term{canonical}{典范}\Term{isomorphism}{同构}
\[\begin{aligned}
\operatorname{R\mathcal Hom}(F^\vee\otimes^{\mathrm L}E,A)&\xrightarrow{\sim}\operatorname{R\mathcal Hom}(F^\vee,E^\vee)&&\text{（7.4）}\\
&\xrightarrow{\sim}F^{\vee\vee}\otimes^{\mathrm L}E^\vee&&\text{（7.6）}\\
&\xrightarrow{\sim}F\otimes^{\mathrm L}E^\vee&&\text{（7.2）。}
\end{aligned}\]
因此，只要核对相容性，即得所需。也可以先\Term{localization}{局部化}再作 dévissage，归结到 $E=F=A$。
\end{proof}

\SourcePage{164}
\textbf{练习 7.9。}设 $E,F\in\operatorname{ob}D(S)_{\underline{\mathrm{parf}}}$。证明存在\Term{canonical}{典范}\Term{isomorphism}{同构}
\[E^\vee\otimes^{\mathrm L}F^\vee\xrightarrow{\sim}(E\otimes^{\mathrm L}F)^\vee.\]

\SGAsection{8}{迹与杯积}
\textbf{8.0。}本节将\Term{finite-type}{有限型}\Term{projective-object}{投射模}的\Term{endomorphism}{自同态}的\Term{trace}{迹}概念推广到\Term{perfect-complex}{完美复形}（应用见 SGA 5 III）。假设与记号同（7.0）。

\textbf{8.1。}设 $E\in\operatorname{ob}D(S)_{\underline{\mathrm{parf}}}$（5.8.4 b））。存在自然\Term{arrow}{箭头}
\[\tau:\operatorname{R\mathcal Hom}(E,E)\longrightarrow A,\]
它由（7.7）的\Term{isomorphism}{同构}的逆 $\operatorname{R\mathcal Hom}(E,E)\to E^\vee\otimes^{\mathrm L}E$ 与（7.5）的\Term{canonical}{典范}\Term{arrow}{箭头} $E^\vee\otimes^{\mathrm L}E\to A$ 复合而成。

设 $U\in\operatorname{ob}\mathscr S$。对 $\tau$ 应用\Term{functor}{函子} $H^0(U,\cdot)$，并注意到 $H^0(U,\operatorname{R\mathcal Hom}(E,E))$ 等同于 $\operatorname{Hom}_{D(U)}(E|U,E|U)$，由此得到 $H^0(U,A)$-\Term{module-homomorphism}{模同态}
\[\operatorname{Tr}_{E|U}:\operatorname{Hom}_{D(U)}(E|U,E|U)\longrightarrow H^0(U,A),\]
称为\Term{trace-homomorphism}{迹同态}。\footnote{事实上，对于 $\operatorname{Tr}_{E|U}$ 的定义，$E$ 具有\Term{finite}{有限}\Term{tor-dimension}{Tor 维数}的假设是多余的；见（8.5）。}

进一步讨论之前，先给出几个直接由定义得到的\Term{trace}{迹}的性质。

\SourcePage{165}
\textbf{8.1.1。}若 $E=A$，则 $\operatorname{Tr}_{A|U}:\operatorname{Hom}(A|U,A|U)\to H^0(U,A)$ 正是熟知的\Term{isomorphism}{同构}。若 $E$ 是一个\Term{finite-free}{有限型自由模}的\Term{direct-summand}{直和因子}，则 $U$ 上方从 $E$ 到 $E$ 的一个\Term{homomorphism}{同态}\Term{local-scheme}{局部}由
\[e\longmapsto\sum_i\langle e,x_i^\vee\rangle x_i\]
给出，其中 $x_i$（相应地，$x_i^\vee$）是 $E$（相应地，$E^\vee$）的\Term{local-scheme}{局部}\Term{section}{截面}；相应的\Term{trace}{迹}为 $\sum_i\langle x_i,x_i^\vee\rangle$。容易由此\Term{pushout}{推出}：若 $E$ 和 $F$ 是\Term{finite-free}{有限型自由模}的\Term{direct-summand}{直和因子}，且 $u:(E\oplus F)|U\to(E\oplus F)|U$ 由矩阵 $\begin{pmatrix}a&b\\c&d\end{pmatrix}$ 定义，则
\[\operatorname{Tr}_{(E\oplus F)|U}(u)=\operatorname{Tr}_{E|U}(a)+\operatorname{Tr}_{F|U}(d).\]

\textbf{8.1.2。}设 $E$ 是一个\Term{strictly-perfect}{严格完美}\Term{complex}{复形}，$U\in\operatorname{ob}\mathscr S$，$u:E|U\to E|U$ 是\Term{complex}{复形}\Term{homomorphism}{同态}，其\Term{set-component}{分量}为 $u^i:E^i|U\to E^i|U$。则
\[\operatorname{Tr}_{E|U}(u)=\sum_i(-1)^i\operatorname{Tr}_{E^i|U}(u^i).\]
事实上，问题是\Term{local-scheme}{局部}的，故可设 $u$ 由 $\mathcal H^0(E^\vee\otimes E)$ 的一个\Term{section}{截面}定义，甚至可设此\Term{section}{截面}形如 $\prod_i\sum_j x_{ij}^\vee\otimes x_{ij}$，其中 $x_{ij}$（相应地，$x_{ij}^\vee$）是 $E^i$（相应地，$E^{i\vee}$）的\Term{section}{截面}。将（7.5）的\Term{arrow}{箭头} $E^\vee\otimes E\to A$ 明确写出，即得
\[\operatorname{Tr}_{E|U}(u)=\sum_{i,j}(-1)^i\langle x_{ij},x_{ij}^\vee\rangle=\sum_i(-1)^i\operatorname{Tr}_{E^i|U}(u^i),\]
证毕。

\textbf{8.2。}设 $E,F\in\operatorname{ob}D(S)_{\underline{\mathrm{parf}}}$。已由（7.8）知存在\Term{canonical-homomorphism}{典范同态}
\[\gamma:\operatorname{R\mathcal Hom}(E,F)\otimes^{\mathrm L}\operatorname{R\mathcal Hom}(F,E)\longrightarrow A.\]
设 $U\in\operatorname{ob}\mathscr S$，并给定 $u\in H^0(U,\operatorname{R\mathcal Hom}(E,F))$ 和 $v\in H^0(U,\operatorname{R\mathcal Hom}(F,E))$。
\SourcePage{166}
可将 $u$（相应地，$v$）解释为 $D(U)$ 中从 $A|U$ 到 $\operatorname{R\mathcal Hom}(E,F)|U$（相应地，$\operatorname{R\mathcal Hom}(F,E)|U$）的一个\Term{homomorphism}{同态}，并构造
\[u\otimes^{\mathrm L}v:A|U\longrightarrow(\operatorname{R\mathcal Hom}(E,F)\otimes^{\mathrm L}\operatorname{R\mathcal Hom}(F,E))|U.\]
映射 $(u,v)\mapsto\langle u,v\rangle=\gamma\circ(u\otimes^{\mathrm L}v)$ 定义 $H^0(U,A)$-\Term{linear}{线性}映射
\[\operatorname{Hom}_{D(U)}(E|U,F|U)\otimes_{H^0(U,A)}\operatorname{Hom}_{D(U)}(F|U,E|U)\longrightarrow H^0(U,A),\]
我们称其为\Term{cup-product}{杯积}。\footnote{与定义 $\operatorname{Tr}_{E|U}$ 时的注记相同。}

\Term{trace}{迹}与\Term{cup-product}{杯积}的关系由下列命题表达。
\begin{SGAresult}{命题}{8.3}
在（8.2）的情形，对于 $u\in\operatorname{Hom}(E|U,F|U)$ 和 $v\in\operatorname{Hom}(F|U,E|U)$，有
\[\operatorname{Tr}_{E|U}(vu)=\operatorname{Tr}_{F|U}(uv)=\langle u,v\rangle.\]
特别地，若 $E=F$，则
\[\operatorname{Tr}_{E|U}(u)=\langle u,\operatorname{Id}\rangle=\langle\operatorname{Id},u\rangle.\]
\end{SGAresult}
\begin{proof}
为简便起见，设 $U=S$。于是有两个\Term{arrow}{箭头} $u:A\to E^\vee\otimes^{\mathrm L}F$ 和 $v:A\to F^\vee\otimes^{\mathrm L}E$，从而有
\[u\otimes^{\mathrm L}v:A\longrightarrow E^\vee\otimes^{\mathrm L}F\otimes^{\mathrm L}F^\vee\otimes^{\mathrm L}E.\]
另一方面，有 $vu:A\to E^\vee\otimes^{\mathrm L}E$ 和 $uv:A\to F^\vee\otimes^{\mathrm L}F$。命题由下列图的\Term{commutative-diagram}{交换性}得到（其核对是“纯粹机械的”）：
\SourcePage{167}
\[\begin{tikzcd}[column sep=large]
A\arrow[r,"u\otimes^{\mathrm L}v"]\arrow[dr,"vu"']&E^\vee\otimes^{\mathrm L}F\otimes^{\mathrm L}F^\vee\otimes^{\mathrm L}E\arrow[d,"\operatorname{Id}\otimes^{\mathrm L}\tau_F"]\\
&E^\vee\otimes^{\mathrm L}E
\end{tikzcd}\]
\[\begin{tikzcd}[column sep=large]
A\arrow[r,"u\otimes^{\mathrm L}v"]\arrow[dr,"uv"']&E^\vee\otimes^{\mathrm L}F\otimes^{\mathrm L}F^\vee\otimes^{\mathrm L}E\arrow[d,"\tau_E\otimes^{\mathrm L}\operatorname{Id}"]\\
&F^\vee\otimes^{\mathrm L}F
\end{tikzcd}\]
\[\begin{tikzcd}[column sep=large]
E^\vee\otimes^{\mathrm L}F\otimes^{\mathrm L}F^\vee\otimes^{\mathrm L}E\arrow[r,"\operatorname{Id}\otimes^{\mathrm L}\tau_E"]\arrow[d,"\operatorname{Id}\otimes^{\mathrm L}\tau_F"']&F^\vee\otimes^{\mathrm L}F\arrow[d,"\tau_F"]\\
E^\vee\otimes^{\mathrm L}E\arrow[r,"\tau_E"']&A.
\end{tikzcd}\]
\end{proof}

\begin{SGAresult}{推论}{8.4}
$E,F,U$ 如（8.2）中所给，设 $u\in\operatorname{Hom}(E|U,E|U)$，$s:E|U\to F|U$ 为\Term{isomorphism}{同构}，从而 $sus^{-1}\in\operatorname{Hom}(F|U,F|U)$。则
\[\operatorname{Tr}_{F|U}(sus^{-1})=\operatorname{Tr}_{E|U}(u).\]
\end{SGAresult}
\begin{proof}
这是立即可见的。
\end{proof}

\textbf{注记 8.5。}若 $\mathcal C$ 是 $\mathscr S$-\Term{fibered-category}{纤维范畴}，以 $\operatorname E(\mathcal C)$ 表示如下 $\mathscr S$-\Term{fibered-category}{纤维范畴}：其在 $U\in\operatorname{ob}\mathscr S$ 处的\Term{fiber}{纤维}是以对 $(L,u)$ 为\Term{object}{对象}的\Term{category}{范畴}，其中 $L\in\operatorname{ob}\mathcal C_U$，$u\in\operatorname{End}_{\mathscr S}(L)$；从 $(L,u)$ 到 $(M,v)$ 的\Term{arrow}{箭头}为满足 $fu=vf$ 的\Term{arrow}{箭头} $f:L\to M$。由（8.4），可将\Term{trace-homomorphism}{迹同态}视为从 $\operatorname E(D(\mathscr S)_{\underline{\mathrm{parf}}})^{\mathrm{is}}$ 到 $\mathcal A$ 的 $\mathscr S$-\Term{functor}{函子}（其中上标 $\mathrm{is}$ 表示仅考虑\Term{isomorphism}{同构}作为\Term{arrow}{箭头}，而 $\mathcal A$ 为由 $A$ 定义的具有\Term{discrete}{离散}\Term{fiber}{纤维}的 $\mathscr S$-\Term{category}{范畴}，参见（6.6 a）））。因为每个\Term{perfect-complex}{完美复形}\Term{local-scheme}{局部}具有\Term{finite}{有限}\Term{tor-dimension}{Tor 维数}，遗忘 $\mathscr S$-\Term{functor}{函子}
\[\operatorname E(D(\mathscr S)_{\underline{\mathrm{parf}}})^{\mathrm{is}}\longrightarrow\operatorname E(D(\mathscr S)_{\mathrm{parf}})^{\mathrm{is}}\]
在\Term{associated-stack}{伴随叠}上诱导\Term{isomorphism}{同构}。因而，
\SourcePage{168}
\Term{trace-homomorphism}{迹同态}唯一延拓为从 $\operatorname E(D(\mathscr S)_{\mathrm{parf}})^{\mathrm{is}}$ 到 $\mathcal A$ 的 $\mathscr S$-\Term{functor}{函子}。对于 $E\in\operatorname{ob}\operatorname{Parf}(U)$（$U\in\operatorname{ob}\mathscr S$），仍以
\[\operatorname{Tr}_E:\operatorname{Hom}(E,E)\longrightarrow H^0(U,A)\]
表示相应\Term{homomorphism}{同态}，它是 $H^0(U,A)$-\Term{module-homomorphism}{模同态}。同样可见，\Term{cup-product}{杯积}的定义（8.2），以及断言（8.3）和（8.4），都可推广到\Term{perfect-complex}{完美复形}的情形（不必具有\Term{finite}{有限}\Term{tor-dimension}{Tor 维数}）。

\begin{SGAresult}{命题}{8.6}
设 $U\in\operatorname{ob}\mathscr S$，$E,F\in\operatorname{Parf}(U)$，$u\in\operatorname{Hom}(E,E)$，$v\in\operatorname{Hom}(F,F)$，从而
\[u\otimes^{\mathrm L}v\in\operatorname{Hom}(E\otimes^{\mathrm L}F,E\otimes^{\mathrm L}F).\]
则
\[\operatorname{Tr}_{E\otimes^{\mathrm L}F}(u\otimes^{\mathrm L}v)=\operatorname{Tr}_E(u)\operatorname{Tr}_F(v).\]
\end{SGAresult}
\begin{proof}
问题是\Term{local-scheme}{局部}的，故由（8.4）与（4.10），可设 $u$ 与 $v$ 为 $K^b(\mathcal C_0)$ 的\Term{arrow}{箭头}。使用（8.1.2），归结到 $u$ 与 $v$ 为 $\mathcal C_0$ 的\Term{arrow}{箭头}的情形。最后借助（8.1.1）即得证。
\end{proof}

\SourcePage{169}
\SGAsection{}{参考文献}
\begin{description}
\item[{[H]}] Hartshorne, R.，《剩余与\Term{duality}{对偶性}》（\textit{Residues and Duality}），Lecture Notes 第 20 号，Springer，1966。
\item[{[V]}] Verdier, J.-L.，《\Term{derived-category}{导出范畴}》（\textit{Catégories dérivées}），I.H.E.S.，1963。
\item[{[G]}] Giraud, J.，《非 Abel \Term{cohomology}{上同调}》（\textit{Cohomologie non abélienne}），即将由 North Holl. Pub. Comp. 出版。
\item[{SGA 4}] 《\Term{scheme}{概形}的\Term{etale-cohomology}{平展上同调}》（\textit{Cohomologie étale des schémas}），M. Artin、A. Grothendieck 与 J.-L. Verdier 著，即将由 North Holl. Pub. Comp. 出版。
\item[{SGA 5}] 《$\ell$-进\Term{cohomology}{上同调}与 $L$ 函数》（\textit{Cohomologie $\ell$-adique et fonctions L}），A. Grothendieck 著，即将由 North Holl. Pub. Comp. 出版。
\item[{[1]}] Hakim, M.，《\Term{ringed-topos}{赋环拓扑斯}与\Term{relative-scheme}{相对概形}》（\textit{Topos annelés et schémas relatifs}），学位论文（撰写中）。
\item[{[2]}] Verdier, J.-L.，《\Term{abelian-category}{Abel 范畴}的\Term{derived-category}{导出范畴}》（\textit{Les catégories dérivées des catégories abéliennes}），即将由 North Holl. Pub. Comp. 出版。
\item[{[MULT]}] Cartan, H. 与 Eilenberg, S.，《\Term{homological-algebra}{同调代数}》（\textit{Homological Algebra}），Princeton，1956。
\end{description}
